\chapter{La simetría}\label{ch:g4:symmetry}

Mariposas, \omterm{def:g2:shapes:solids}{caras}, copos de nieve: a la naturaleza le gustan las figuras cuyas dos \omterm{def:g3:division:half}{mitades}
encajan al doblarlas. Este capítulo explora la simetría con dobleces y papel
cuadriculado --- las construcciones con la escuadra esperan al
\cref{ch:g6:symmetry}.

\section{Doblar y encajar}

\begin{definition}[Eje de simetría]\label{def:g4:symmetry:def}
Una recta es un \emph{eje de simetría}\index{eje de simetría} de una
figura cuando, al doblar la figura por esa recta, las dos \omterm{def:g3:division:half}{mitades}
encajan \emph{exactamente}. Una figura puede no tener ningún eje, tener uno o
tener varios.
\end{definition}

\begin{omfigure}
\begin{tikzpicture}[scale=0.6]
  % figura simetrica: triangulo isosceles
  \draw[thick, fill=omDef!12] (0,0) -- (2.4,0) -- (1.2,2.6) -- cycle;
  \draw[omThm, dashed, thick] (1.2,-0.5) -- (1.2,3.0);
  \node[below, font=\small] at (1.2,-0.7) {un eje};
  % rectangulo: dos ejes
  \begin{scope}[xshift=5.4cm]
  \draw[thick, fill=omDef!12] (0,0.3) rectangle (3.2,2.3);
  \draw[omThm, dashed, thick] (1.6,-0.4) -- (1.6,3.0);
  \draw[omThm, dashed, thick] (-0.5,1.3) -- (3.7,1.3);
  \node[below, font=\small] at (1.6,-0.7) {dos ejes};
  \end{scope}
  % escaleno: ninguno
  \begin{scope}[xshift=11.2cm]
  \draw[thick, fill=omDef!12] (0,0) -- (3.0,0.4) -- (0.8,2.2) -- cycle;
  \node[below, font=\small] at (1.4,-0.7) {ningún eje};
  \end{scope}
\end{tikzpicture}

{\small Dobla por la línea discontinua: las \omterm{def:g3:division:half}{mitades} encajan. El último
triángulo no tiene ningún doblez que funcione --- ningún \omterm{def:g4:symmetry:def}{eje de simetría}.}
\end{omfigure}

\begin{example}[Comprobar con papel de calco]\label{ex:g4:symmetry:trace}
Calca la figura, dale la vuelta al papel de calco por el eje
candidato y vuelve a ponerlo: si el calco cubre la figura exactamente, el
eje es auténtico. La vista sola engaña a menudo --- la diagonal de un rectángulo
\emph{parece} un eje, pero la prueba del doblez dice que no (¡pruébalo!).
\end{example}

\section{Completar figuras en la cuadrícula}

\begin{method}[Simétrica de una figura respecto de una línea de la cuadrícula]\label{met:g4:symmetry:grid}
El eje es una línea vertical (u horizontal) de la cuadrícula.
\begin{enumerate}
  \item toma uno a uno cada \omterm{def:g4:geometry:polygon}{vértice} de la figura;
  \item cuenta su distancia al eje en casillas;
  \item coloca el \omterm{def:g4:geometry:polygon}{vértice} imagen a la \emph{misma} distancia al
        \emph{otro} lado, en la misma fila (o columna);
  \item une los \omterm{def:g4:geometry:polygon}{vértices} imagen en el mismo orden y comprueba: la
        imagen es la gemela reflejada de la figura.
\end{enumerate}
\end{method}

\begin{omfigure}
\begin{tikzpicture}[scale=0.55]
  \draw[gray!40, thin] (0,0) grid (12,5);
  \draw[omThm, very thick] (6,-0.4) -- (6,5.4);
  \draw[thick, omDef, fill=omDef!15]
    (1,1) -- (5,1) -- (5,2) -- (3,2) -- (3,4) -- (1,4) -- cycle;
  \draw[thick, omMeth, fill=omMeth!15]
    (11,1) -- (7,1) -- (7,2) -- (9,2) -- (9,4) -- (11,4) -- cycle;
  \draw[dashed] (5,2) -- (7,2);
  \draw[dashed] (3,4) -- (9,4);
\end{tikzpicture}

{\small Completar respecto del eje rojo: cada \omterm{def:g4:geometry:polygon}{vértice} salta a la misma
distancia al otro lado. Los \omterm{def:g4:geometry:polygon}{vértices} que tocan el eje se quedarían
donde están.}
\end{omfigure}

\begin{example}[La misma forma, el mismo tamaño, del revés]\label{ex:g4:symmetry:properties}
La imagen tiene exactamente las mismas longitudes y los mismos ángulos que el
original: la simetría copia la figura --- pero le da la vuelta, como una
mano izquierda y una mano derecha. Si la letra original es una F, su imagen
reflejada es una F al revés, no otra F.
\end{example}

\section{Los ejes de las figuras habituales}

\begin{example}[Contar ejes]\label{ex:g4:symmetry:shapes}
Doblando:
\begin{itemize}
  \item un cuadrado: $4$ ejes (dos por los puntos medios de lados
        opuestos y dos por las diagonales);
  \item un rectángulo: $2$ ejes (¡solo los de los puntos medios de lados opuestos!);
  \item un triángulo isósceles: $1$; un triángulo equilátero: $3$;
  \item un círculo: cualquier recta que pase por su centro --- más ejes de los
        que se pueden contar.
\end{itemize}
\end{example}

\begin{example}[La simetría en el alfabeto]\label{ex:g4:symmetry:letters}
A, M, T, U y V tienen eje vertical; B, C, D y E tienen uno
horizontal; H, I, O y X tienen los dos. Las palabras hechas con las letras adecuadas pueden
ser simétricas también: OSO tiene eje vertical; OXO funciona en los dos sentidos.
\end{example}

\section{Ejercicios}

\begin{exercise}[$\star$]\label{exo:g4:symmetry:1}
Calca y dobla: ¿cuáles de estas figuras tienen un \omterm{def:g4:symmetry:def}{eje de simetría} --- un
cuadrado; un rectángulo (no cuadrado); un triángulo escaleno (con los tres lados
distintos); un círculo?
\end{exercise}

\begin{exercise}[$\star$]\label{exo:g4:symmetry:2}
Traza los \omterm{def:g4:symmetry:def}{ejes de simetría} de: un cuadrado; un rectángulo; un triángulo
equilátero. ¿Cuántos tiene cada uno?
\end{exercise}

\begin{exercise}[$\star$]\label{exo:g4:symmetry:3}
Usa la prueba del papel de calco del \cref{ex:g4:symmetry:trace} para demostrar
que la diagonal de un rectángulo de $6 \times 4$ \emph{no} es un \omterm{def:g4:symmetry:def}{eje
de simetría}.
\end{exercise}

\begin{exercise}[$\star$]\label{exo:g4:symmetry:4}
En papel cuadriculado, traza un eje vertical y dibuja la letra L (de tres casillas
de alto y dos de ancho) a $2$ casillas del eje. Construye su imagen
reflejada.
\end{exercise}

\begin{exercise}[$\star$]\label{exo:g4:symmetry:5}
En papel cuadriculado, traza un eje horizontal y dibuja encima una banderita (un triángulo
de $1 \times 1$ sobre un mástil de $3$ casillas). Completa la figura por
simetría por debajo del eje.
\end{exercise}

\begin{exercise}[$\star$]\label{exo:g4:symmetry:6}
Una figura toca su \omterm{def:g4:symmetry:def}{eje de simetría} en un punto $P$. ¿Dónde está la
imagen de $P$? Explícalo con el dibujo del doblez.
\end{exercise}

\begin{exercise}[$\star$]\label{exo:g4:symmetry:7}
Clasifica las letras mayúsculas A, B, E, F, H, N, O y T: ¿eje vertical?
¿eje horizontal? ¿los dos? ¿ninguno?
\end{exercise}

\begin{exercise}[$\star$]\label{exo:g4:symmetry:8}
Un triángulo tiene lados de $5$ cm, $5$ cm y $3$ cm. ¿Tiene algún
\omterm{def:g4:symmetry:def}{eje de simetría}? ¿Qué lados intercambia el doblez?
\end{exercise}

\begin{exercise}[$\star\star$]\label{exo:g4:symmetry:9}
En papel cuadriculado, dibuja una figura tuya que tenga \emph{exactamente dos}
\omterm{def:g4:symmetry:def}{ejes de simetría} y traza los dos. (Pista: piensa en un rectángulo
o en una cruz.)
\end{exercise}

\begin{exercise}[$\star\star$]\label{exo:g4:symmetry:10}
Se da la \omterm{def:g3:division:half}{mitad} de un dibujo simétrico: la \omterm{def:g3:division:half}{mitad} izquierda de una casa
(una pared cuadrada y medio tejado triangular), con un eje vertical
a lo largo de su borde derecho. Describe o dibuja la casa completa. ¿Cómo son
las longitudes de la \omterm{def:g3:division:half}{mitad} derecha completada frente a las de la izquierda?
\end{exercise}

\begin{exercise}[$\star\star$]\label{exo:g4:symmetry:11}
Aya afirma: «todo triángulo con dos lados iguales tiene un \omterm{def:g4:symmetry:def}{eje de
simetría}, y todo triángulo con un \omterm{def:g4:symmetry:def}{eje de simetría} tiene dos lados
iguales». Pon a prueba su doble afirmación con dibujos: un triángulo isósceles,
uno equilátero y uno escaleno. ¿La respalda el doblez?
\end{exercise}
